% 19-limitations.tex: нынешние ограничения и то, какая внешняя проверка требуется.
\section{Нынешние ограничения и внешняя проверка, которая всё ещё требуется}
\label{sec:limitations}

Книга готовности репозитория открывается одной строкой, и настоящий документ повторяет её как
собственную границу: на \code{1.0.0-кв.7} Полярис не готов к работе с настоящими
идентификационными данными. Каждый инженерный пробел, перечисленный в книге, закрыт и закреплён
проверкой; то, что остаётся, построить внутри репозитория нельзя. Перечень ниже отделяет то, что
не сделано, от того, чего автор сделать не может, и в одном отношении он длиннее, чем в версии 2,
потому что в прошедший промежуток репозиторий находил пределы, а не устранял их.

\subsection{Решения, которые может принять лишь развёртывающая организация}

Прежде чем можно будет держать настоящие данные, для названного развёртывания должны быть
записаны девять решений: правовое основание, оценка воздействия и одобрение регулятора; драйвер
хранения подписывающего ключа и то, кто этот ключ держит; топология отказоустойчивости и её хосты
и то, будет ли работать второй регион; шифрование хранимых данных и хранитель его ключа; место и
расписание удалённого резервного копирования; подсистема оповещений, график дежурств и то, кто
получает вызов о принуждении; политика права на стирание; расписание хранения для каждого класса
данных вместе с юристом, который подтверждает, что число удовлетворяет закону; и независимая
проверка на проникновение с утверждением модели угроз. Полярис записывает каждое решение и его
обоснование; закона он не знает. \sExt

\subsection{Что не построено}

\begin{itemize}
  \item Физический токен есть спецификация, эмулятор и опубликованные векторы; ни одна карта не
    произведена и ни один кристалл не сертифицирован, а профиль называет те свои свойства, которых
    эмулятор установить не способен.
  \item Ни один аппаратный модуль безопасности не применялся; драйвер ПККС\#11 проверен против
    программного токена, драйвер СУК против заместителя, а решёточная библиотека не прошла
    проверку по ФИПС.
  \item Второй регион есть резервный кластер с ненулевой точкой восстановления; раздача состояний
    есть набор правил кэширования, а не развёртывание; а национальная нагрузка смоделирована в
    масштабе и никогда не запускалась.
  \item Хеш-основанный подписант зарегистрирован, но не подключён. Переход от пулера к базе
    данных, подписи сертификатов и нынешние аутентификаторы операторов остаются классическими.
  \item Верификатор не способен узнать, что алгоритм объявлен устаревшим, поэтому вывод
    алгоритма из обращения есть выпуск программного обновления для каждой доверяющей стороны;
    откат и смешанное окно частичной миграции не измерены.
  \item Ослеплённые и одноразовые предъявления не построены, поэтому обыкновенное удостоверение
    показывает верификатору устойчивое значение токена, подпись эмитента и ключ держателя, а
    предъявление СД-ЖВТ ПУ показывает девять подписанных эмитентом значений, одинаковых повсюду.
    Дескриптор, свой для каждой доверяющей стороны, ограничивает то, что верификатор
    \emph{хранит}; лишь путь с нулевым разглашением ограничивает то, что ему \emph{показывают}, и
    эта граница есть размер эпохи, которому ничто не задаёт нижнего предела выше единицы и о
    котором не сообщает ни один вердикт.
  \item Ни нуллификатор с областью, ни попарный дескриптор не прячут держателя от
    \emph{эмитента}, который выводит каждый лист эпохи, чтобы построить дерево. Вывод, слепой для
    эмитента, не заявляется.
  \item Как эмитент Полярис не имеет удостоверения с признаками и анонимного условия; как
    верификатор он принимает раскрытое логическое значение, которое эмитенту пришлось
    предусмотреть.
  \item Механизм принуждения противостоит беспечному принуждающему и не противостоит ни
    осведомлённому принуждающему рядом с держателем, ни законному или ведомственному доступу,
    против которого его неизгладимая запись в итоге вредит. Сигнала тревоги со стороны держателя
    нет, и он не предлагается.
  \item Подлинность документов при регистрации принимается со слов оператора; уровень
    достоверности выводится из того, какие свидетельства были предъявлены, а не из того, были ли
    они подлинными.
  \item Из полей идентификаторов, растущих с темпом регистрации, два остаются 32-битными,
    \code{ФизическоеЛицо.физ\_лицо\_ид} и \code{ТокенЛичности.токен\_ид}: запас примерно на двадцать девять
    лет при заявленном всплеске в 200\,000 в день, за горизонтом модели ёмкости; временное удостоверение на период выдержки при восстановлении не
    построено; вердикт о закате опирается на заверения оператора, которых система проверить не
    может.
\end{itemize}

\subsection{Что сделал кто-то другой, а что нет}

\begin{itemize}
  \item Один неизменённый кошелёк, написанный кем-то другим, предъявил удостоверение и был принят,
    на одном пути, в одном формате, с одним алгоритмом подписи, а перед тем один раз Полярис
    отверг. Ни второго кошелька, ни кошелька на родном протоколе, ни сторонней доверяющей
    стороны.
  \item Размещённый пакет Фонда ОпенАйДи прогнал профиль верификатора: одиннадцать модулей, ноль
    отказов, семь отрицательных модулей засчитаны собственной оценкой службы, четыре положительных
    ждут на рассмотрении человека, которого Фонд назначает при подаче. Не подано, не
    сертифицировано.
  \item Опубликованный сторонний набор векторов сходится с обоими свидетелями подписи на
    примитиве МЛ-ДСА-65. Он проверяет примитив, а не протокол.
  \item Четыре пакета лежат в публичных реестрах, и каждый проверен установкой из реестра в
    чистое окружение. Никто за пределами репозитория не сообщал, что установил хоть один, а число
    скачиваний в реестре не есть человек.
  \item Ни внешней проверки на проникновение, ни криптографического аудита, ни работы красной
    команды, ни программы вознаграждений за уязвимости, ни сертификации не было. Документ об
    объёме работ и пакет для рецензента существуют; заказа не было. Почти всякая гарантия в
    настоящем документе проверяется механизмами, написанными тем же автором, что и сама
    гарантия.
  \item Ни один оператор, кроме автора, ничего из этого не запускал, и это единственное условие,
    отделяющее кандидата в выпуски от \code{1.0.0}. Ни одного пилота с согласившимися
    пользователями не проводилось; данные всегда были условными.
  \item Ни одна сторонняя команда не реализовала родной протокол по одной лишь спецификации и
    набору проверок соответствия; набор есть договор, а интеграции не случилось.
  \item Ни один институционально независимый свидетель не наблюдает ни за одним журналом;
    свидетели в записи суть процессы в учении. Два федеративных экземпляра суть два экземпляра
    одной кодовой базы, запущенные одним автором.
  \item Схема поверх библиотеки доказательств не проходила внешнего ревью, и конкретная стойкость в
    битах поставляемой конфигурации доказательств независимо не выводилась.
  \item Тезис репозитория, что посторонний способен прочесть его с нуля и правильно вывести то, что
    ему нужно, до своего срока ни разу не проверялся посторонним и записан как неопределённый, а
    сильное утверждение навсегда снято; путь постороннего из раздела~\ref{sec:use} есть более узкая
    проверка, которая его заменила, и никто извне пока не сообщал, что прошёл его.
\end{itemize}

\subsection{Что представляют собой числа}

Каждое число о производительности в настоящем документе относится к одному узлу на одном
ноутбуке, либо к моделированию в указанном масштабе, либо к проекции, которую репозиторий так и
помечает, называя допущение. Ни один промышленный кластер не измерялся, а единственное измерение
на настоящей топологии переключения есть медиана при двух запросах в секунду, причём хвост
распределения не приводится, потому что выборка его не выдерживает. Национальные плановые цели
суть цели; арифметика, которая их перекрывает, есть лёгкая половина вопроса, а половина, которая
едва не отказала, была целочисленным полем.

\subsection{Документация, отстающая от кода}

Работа над настоящей версией нашла места, где собственная проза репозитория отстаёт от его кода
или истории; они записаны в авторских заметках, сопровождающих эту версию, а не исправлены
молча, как это было с версией 2. Там, где настоящий документ расходится с каким-либо документом
репозитория, документ следовал коду на \code{1.0.0-кв.7} и говорит об этом в заметке.
